
\begin{proposition}[Counterexample: A specific implication in the printed proof of Proposition 7.3]\label{cdp:st:A.4.2}
The implication that a ruled minimal surface $S_0$ with a one-parameter family of numerically trivial line bundles $M_t$ and nonzero sections of $\Omega^1_{S_0}\otimes M_t$ must be ruled over a curve of genus at least two is false when the family $M_t$ is allowed to be constant, as the printed proof explicitly allows.
\end{proposition}
\begin{proof}
Take $S_0=E\times\PP^1$ for an elliptic curve $E$ and take $M_t=\OO_{S_0}$ for every parameter. The product cotangent decomposition gives
\[
 h^q(S_0,\Omega^1_{S_0})=(1,2,1),\qquad q=0,1,2.
\]
Indeed the $E$ cotangent summand contributes $(1,1,0)$, while the $\PP^1$ cotangent summand contributes $(0,1,1)$ by the Kunneth formula and $h^1(\PP^1,\OO(-2))=1$. A nonzero invariant form on $E$ provides the asserted sections for every parameter, whereas the ruling has genus one. This refutes the stated implication. It is a counterexample to that implication with its explicit data; it supplies no threefold satisfying all of Proposition~7.3's ambient hypotheses.
\end{proof}


